\section{Atomic, Molecular \& Optical}
\subsection[Challenge 4: Orbital angular momentum conservation in high harmonic generation]{Challenge 4: Orbital angular momentum conservation in high harmonic generation}
\textbf{Problem ID:} \texttt{Challenge\_4\_main}\\
\textbf{Primary inferred discipline:} Atomic, Molecular \& Optical\\
\textbf{Secondary inferred disciplines:} Nonlinear Dynamics\\
\textbf{Python implementation:} \texttt{challenges/04/answer.py}

\subsubsection*{Problem}
\paragraph*{Problem setup}
High-harmonic generation (HHG) is a nonlinear optical process in which intense laser fields interact with gas-phase atoms or molecules to produce coherent extreme ultraviolet (EUV) radiation. In this setup, a sequence of time-delayed laser pulses is focused into a gas jet to investigate the role of structured light in controlling the HHG process.

\paragraph*{Main problem}

I create a driving field composed of three temporally identical laser pulses of 800 nm center wavelength, 50 fs FWHM pulse duration, with the peaks of the pulses at $t = 0$ fs, $t = 30$ fs, and $t = 60$ fs, respectively. The first pulse is left circularly polarized and carries an orbital angular momentum (OAM) of $\ell = -1$, the second pulse is right circularly polarized and has OAM $\ell = 2$, and the third pulse is left circularly polarized and has OAM $\ell = 1$. The pulses are focused, overlapped in space, into a gas jet, and used to drive high-harmonic generation to produce EUV light. Using the convention $\sigma = +1$ for left circular polarization and $\sigma = -1$ for right circular polarization, calculate the possible OAM and helicity states, reporting as a set of $(\ell, \sigma)$ pairs, that appear in the 23rd harmonic order.
